\chapter{Standard Form, Basic Solutions, and the Matrix View}

The simplex method does not work on inequalities. It walks from corner to
corner of the feasible region, and at each corner it solves a system of
linear \emph{equations}. This chapter builds the bridge from the modelling
form of an LP to that computational form: every LP can be rewritten in
\emph{standard form} (a maximisation with equality
constraints and nonnegative variables) by introducing slack and surplus
variables, splitting free variables, and flipping minimisations.

Once the LP reads $Ax=b$, $x\ge 0$, its corners acquire an algebraic
fingerprint: the \emph{basic feasible solutions}. Choosing which variables
are ``basic'' and setting the rest to zero produces at most
$\binom{n}{m}$ candidate points, and the corners of the feasible region are
exactly the feasible ones among them. Everything the simplex method does in
later chapters is bookkeeping on this correspondence, so the hand skills
drilled here (converting, enumerating, checking feasibility, matching
corners) are used in every problem that follows.

\section{Core Rules}

\begin{rulebox}{The two matrix forms of an LP}
With $A$ an $m\times n$ matrix, $b\in\R^{m}$ and $c,x\in\R^{n}$, the
\emph{inequality form} of an LP is
\[
  \max\; z=c^{T}x
  \quad\text{s.t.}\quad Ax\le b,\;\; x\ge 0 ,
\]
where $Ax\le b$ is read row by row. The \emph{standard form} replaces the inequalities by equations:
\[
  \max\; z=c^{T}x
  \quad\text{s.t.}\quad Ax=b,\;\; x\ge 0 .
\]
The simplex method operates exclusively on the standard form.
\end{rulebox}

\begin{rulebox}{Slack and surplus variables}
Each inequality is turned into an equation by one new nonnegative variable
with objective coefficient $0$:
\begin{itemize}[leftmargin=1.6em,itemsep=0.3ex]
  \item \textbf{$\le$ constraint:} $a_{1}x_{1}+\dots+a_{n}x_{n}\le \beta$
        becomes $a_{1}x_{1}+\dots+a_{n}x_{n}+s=\beta$, $s\ge 0$. The
        \emph{slack variable} $s$ measures the unused
        capacity of the constraint.
  \item \textbf{$\ge$ constraint:} $a_{1}x_{1}+\dots+a_{n}x_{n}\ge \beta$
        becomes $a_{1}x_{1}+\dots+a_{n}x_{n}-s=\beta$, $s\ge 0$. Here $s$
        is a \emph{surplus variable}; it is
        \emph{subtracted}.
\end{itemize}
Use a fresh variable for every constraint, never recycle one slack for
two rows.
\end{rulebox}

\begin{rulebox}{Equality constraints and negative right-hand sides}
An equality constraint is already in the required shape: copy it unchanged,
\emph{no} new variable is introduced. If a right-hand side is negative,
multiply the whole row by $-1$ first (which flips $\le$ into $\ge$ and vice
versa), then convert as usual; the simplex start-up in later chapters
assumes $b\ge 0$.
\end{rulebox}

\begin{rulebox}{Free and nonpositive variables}
Standard form requires every variable to be nonnegative.
\begin{itemize}[leftmargin=1.6em,itemsep=0.3ex]
  \item \textbf{Free variable}: if $x_{j}$ has no sign
        restriction, substitute
        \[
          x_{j}=x_{j}^{+}-x_{j}^{-},\qquad x_{j}^{+},x_{j}^{-}\ge 0
        \]
        everywhere ($x_{j}^{+}$ carries the positive part, $x_{j}^{-}$ the
        negative part).
  \item \textbf{Nonpositive variable:} if $x_{j}\le 0$, substitute
        $y_{j}=-x_{j}\ge 0$ everywhere and drop $x_{j}$.
\end{itemize}
After solving, translate back: $x_{j}=x_{j}^{+}-x_{j}^{-}$ resp.\
$x_{j}=-y_{j}$.
\end{rulebox}

\begin{rulebox}{Minimisation problems}
Minimising $c^{T}x$ is the same task as maximising $-c^{T}x$ over the same
feasible set:
\[
  \min\; c^{T}x \;=\; -\max\;\bigl(-c^{T}x\bigr) .
\]
The optimal \emph{points} coincide; only the optimal \emph{value} changes
sign. So: negate the objective, solve the maximisation, and negate the
optimal value at the very end.
\end{rulebox}

\begin{rulebox}{Basic and nonbasic variables, basic solutions}
Let $Ax=b$ have $m$ equations and $n>m$ variables, with the rows of $A$
linearly independent. Choose $m$ columns of $A$ that are linearly
independent; they form a \emph{basis} $B$. The $m$ variables
belonging to these columns are the \emph{basic variables}; the remaining $n-m$ variables are \emph{nonbasic}. Setting all nonbasic variables to $0$ and
solving the resulting $m\times m$ system $Bx_{B}=b$ gives the unique
\emph{basic solution} of that basis. If in addition
$x_{B}\ge 0$, it is a \emph{basic feasible solution}, short BFS.
\end{rulebox}

\begin{rulebox}{Basic feasible solutions are corners}
For the standard form $Ax=b$, $x\ge 0$:
\[
  x \text{ is a corner of the feasible region}
  \iff
  x \text{ is a basic feasible solution.}
\]
Moreover, if the LP has an optimal solution at all, then it has an optimal
\emph{basic feasible} solution. Hence an optimum can be found by comparing
the objective values of the finitely many BFS: which is exactly what
the simplex method organises efficiently.
\end{rulebox}

\begin{rulebox}{Counting candidate bases}
There are at most
\[
  \binom{n}{m}=\frac{n!}{m!\,(n-m)!}
\]
ways to choose $m$ basic variables out of $n$. The true number of corners
is usually smaller, because a chosen column set may be linearly dependent
(no basic solution), the basic solution may have negative entries
(infeasible), or several bases may describe the same corner (degeneracy).
\end{rulebox}

\begin{rulebox}{Degeneracy}
A basic solution is \emph{degenerate} if at least one
\emph{basic} variable equals $0$: the solution then has fewer than $m$
nonzero entries. Geometrically, more constraints are active at the corner
than needed to define it (in the plane: three or more constraint lines
through one point), and several different bases represent the same corner.
Degeneracy is harmless for the theory but will require care in the simplex
method.
\end{rulebox}

\section{Worked Example}

\begin{worked}{Put the running example in standard form, list every basic solution, and match them to the corners of \chref{linprog}{2}}
\wstep{Step 1. Turn each inequality into an equation with a slack variable.}
\[
  2x_{1}+x_{2}\le100
  \;\longrightarrow\;
  2x_{1}+x_{2}+s_{1}=100,\qquad s_{1}\ge0 ,
\]
\[
  x_{1}+x_{2}\le80
  \;\longrightarrow\;
  x_{1}+x_{2}+s_{2}=80,\qquad s_{2}\ge0 .
\]
The slack has a meaning: $s_{1}$ is \emph{unused machine hours}, $s_{2}$ is
\emph{unused labour hours}. Keeping that in mind makes every later tableau
readable.
\wtrap{A ``$\le$'' constraint takes a slack \emph{added}; a ``$\ge$''
constraint takes a surplus \emph{subtracted}, $ax\ge b\to ax-s=b$. Adding in
both cases is the standard error, and it silently reverses the constraint.
Slack and surplus variables are both required to be $\ge0$.}

\wstep{Step 2. The system in matrix form.}
\[
  A=\begin{pmatrix}2&1&1&0\\1&1&0&1\end{pmatrix},
  \qquad
  \mathbf x=\begin{pmatrix}x_{1}\\x_{2}\\s_{1}\\s_{2}\end{pmatrix},
  \qquad
  \mathbf b=\begin{pmatrix}100\\80\end{pmatrix},
\]
\[
  \max\ \mathbf c^{\!\top}\mathbf x
  \quad\text{s.t.}\quad
  A\mathbf x=\mathbf b,\ \mathbf x\ge\mathbf 0,
  \qquad
  \mathbf c=(40,30,0,0)^{\!\top}.
\]
Slack variables enter the objective with coefficient $0$: idle capacity
earns nothing.

\wstep{Step 3. Count the candidate bases.}
$A$ is $2\times4$ with rank $2$, so a basis is any $2$ of the $4$ columns
that are linearly independent:
\[
  \binom{4}{2}=6\ \text{candidate bases.}
\]
Set the other two variables (the \textbf{non-basic} ones) to zero and solve
for the rest --- that is a \textbf{basic solution}. It is \emph{feasible}
only if every basic variable also comes out $\ge0$.

\wstep{Step 4. Enumerate all six.}
\begin{center}\footnotesize
\begin{tabular}{llll}
  \hline
  basis & solution $(x_{1},x_{2},s_{1},s_{2})$ & feasible? & corner\\
  \hline
  $\{s_{1},s_{2}\}$ & $(0,0,100,80)$ & yes & $(0,0)$\\
  $\{x_{1},s_{2}\}$ & $(50,0,0,30)$ & yes & $(50,0)$\\
  $\{x_{1},x_{2}\}$ & $(20,60,0,0)$ & yes & $(20,60)$\\
  $\{x_{2},s_{1}\}$ & $(0,80,20,0)$ & yes & $(0,80)$\\
  $\{x_{1},s_{1}\}$ & $(80,0,-60,0)$ & \textbf{no} & ---\\
  $\{x_{2},s_{2}\}$ & $(0,100,0,-20)$ & \textbf{no} & ---\\
  \hline
\end{tabular}
\end{center}
Two sample computations. For $\{x_{1},s_{1}\}$: setting $x_{2}=s_{2}=0$, the
labour row gives $x_{1}=80$, and the machine row then gives
$s_{1}=100-160=-60<0$. For $\{x_{2},s_{2}\}$: the machine row gives
$x_{2}=100$, and labour gives $s_{2}=80-100=-20<0$.

\wstep{Step 5. Read the correspondence.}
\[
  \boxed{\text{basic \emph{feasible} solutions}\ \longleftrightarrow\ \text{corners of the feasible region.}}
\]
Four of the six bases are feasible, and the graph in \chref{linprog}{2} had exactly
four corners --- the same four points. The infeasible pair are the
intersections $(80,0)$ and $(0,100)$ that the graph showed lying outside the
region.
\wtrap{$\binom{n}{m}$ counts \emph{candidates}, not corners. Some bases are
infeasible (negative variables) and some column pairs are not even
independent, in which case no basic solution exists at all. The count is an
upper bound, but it does show why enumeration is impractical at scale: a
modest $50\times100$ problem has about $10^{29}$ candidate bases, which is
the entire motivation for the simplex method.}

\wstep{Step 6. Where the optimum must be.}
Every corner is a basic feasible solution and vice versa, and the corner
point theorem puts an optimum at a corner. So the simplex method can restrict
itself to basic feasible solutions --- moving from one to an
\emph{adjacent} better one --- and never look at the interior. That is
\chref{linprog}{4}.

\wstep{Step 7. Degeneracy, defined here so it is recognisable later.}
A basic feasible solution is \textbf{degenerate} if some \emph{basic}
variable equals zero. None of the four above is: at $(20,60)$ the basic
variables $x_{1},x_{2}$ are $20$ and $60$, both positive. Degeneracy means
more than $m$ constraints pass through one corner, so several different
bases describe the same point --- and the simplex method can then pivot
without moving. \chref{linprog}{6} shows what that looks like in a tableau.
\end{worked}

\section{Practice Problems}

Convert by hand, enumerate by hand, and keep the bookkeeping tidy: write
down which variables are nonbasic (set to $0$) before solving for the
basic ones.

\begin{problems}
  \item Convert to standard form:
        \[
          \max\; z=3x_{1}+2x_{2}
          \quad\text{s.t.}\quad
          x_{1}+x_{2}\le 4,\;\;
          x_{1}+3x_{2}\le 6,\;\;
          x_{1},x_{2}\ge 0 .
        \]
  \item Convert to standard form:
        \[
          \max\; z=2x_{1}+5x_{2}
          \quad\text{s.t.}\quad
          x_{1}+2x_{2}\le 8,\;\;
          3x_{1}+x_{2}\ge 3,\;\;
          x_{1},x_{2}\ge 0 .
        \]
  \item Convert to a maximisation in standard form and state how the
        optimal value of the original problem is recovered:
        \[
          \min\; z=4x_{1}+3x_{2}
          \quad\text{s.t.}\quad
          2x_{1}+x_{2}\ge 10,\;\;
          x_{1}+3x_{2}\ge 15,\;\;
          x_{1},x_{2}\ge 0 .
        \]
  \item Convert to standard form:
        \[
          \max\; z=x_{1}+4x_{2}
          \quad\text{s.t.}\quad
          x_{1}+x_{2}=5,\;\;
          x_{1}+2x_{2}\le 8,\;\;
          x_{1},x_{2}\ge 0 .
        \]
  \item Convert to standard form ($x_{2}$ has no sign restriction):
        \[
          \max\; z=2x_{1}+x_{2}
          \quad\text{s.t.}\quad
          x_{1}+x_{2}\le 3,\;\;
          x_{1}-x_{2}\le 2,\;\;
          x_{1}\ge 0,\; x_{2}\ \text{free}.
        \]
  \item Convert to standard form (note the sign restriction on $x_{2}$):
        \[
          \max\; z=3x_{1}-2x_{2}
          \quad\text{s.t.}\quad
          x_{1}-x_{2}\le 6,\;\;
          2x_{1}+x_{2}\le 4,\;\;
          x_{1}\ge 0,\; x_{2}\le 0 .
        \]
  \item Write the following LP in the matrix form $\max\,c^{T}x$ s.t.\
        $Ax\le b$, $x\ge 0$, stating $c$, $A$, $b$ and their dimensions:
        \[
          \max\; z=5x_{1}+4x_{2}+3x_{3}
          \quad\text{s.t.}\quad
          \begin{aligned}[t]
            2x_{1}+3x_{2}+x_{3}&\le 5,\\
            4x_{1}+x_{2}+2x_{3}&\le 11,\\
            3x_{1}+4x_{2}+2x_{3}&\le 8,
          \end{aligned}
          \qquad x_{1},x_{2},x_{3}\ge 0 .
        \]
  \item Bring the LP
        \[
          \max\; z=4x_{1}+3x_{2}
          \quad\text{s.t.}\quad
          2x_{1}+x_{2}\le 10,\;\;
          x_{1}+3x_{2}\le 15,\;\;
          x_{1},x_{2}\ge 0
        \]
        into standard form and write the standard form as a matrix system
        $\bar{A}\bar{x}=b$, $\bar{x}\ge 0$: give $\bar{A}$, $\bar{x}$, $b$
        and the extended cost vector $\bar{c}$.
  \item Find \emph{all} basic solutions of the system
        \[
          \begin{aligned}
            x_{1}+x_{2}+x_{3} &= 4\\
            x_{2}+x_{4} &= 2
          \end{aligned}
          \qquad x_{1},\dots,x_{4}\ge 0 ,
        \]
        by going through all choices of two basic variables. State for each
        choice whether it yields a basis, and whether the basic solution is
        feasible.
  \item Consider the system
        \[
          \begin{aligned}
            x_{1}+x_{2}+x_{3} &= 4\\
            2x_{1}+x_{2}+x_{4} &= 6
          \end{aligned}
          \qquad x_{1},\dots,x_{4}\ge 0 .
        \]
        For each of the following vectors $(x_{1},x_{2},x_{3},x_{4})$
        decide: is it a solution of the equations, is it feasible, and is
        it basic?
        \begin{itemize}[leftmargin=1.8em,itemsep=0.2ex]
          \item[(a)] $(2,2,0,0)$
          \item[(b)] $(1,2,1,2)$
          \item[(c)] $(3,0,1,0)$
          \item[(d)] $(0,6,-2,0)$
          \item[(e)] $(1,1,1,1)$
        \end{itemize}
  \item Bring the LP of Problem~1,
        \[
          \max\; z=3x_{1}+2x_{2}
          \quad\text{s.t.}\quad
          x_{1}+x_{2}\le 4,\;\;
          x_{1}+3x_{2}\le 6,\;\;
          x_{1},x_{2}\ge 0 ,
        \]
        into standard form and enumerate all basic solutions of the system.
        Mark the basic \emph{feasible} solutions.
  \item (Continuation of Problem~11.) Compute all corner points of the
        feasible region in the $(x_{1},x_{2})$-plane by intersecting
        constraint lines by hand. Match every BFS from Problem~11 to a
        corner, interpret the two infeasible basic solutions geometrically,
        and read off the optimum by comparing objective values.
  \item For the LP
        \[
          \max\; z=2x_{1}+3x_{2}
          \quad\text{s.t.}\quad
          x_{1}+2x_{2}\le 6,\;\;
          x_{1}\le 4,\;\;
          x_{1},x_{2}\ge 0 ,
        \]
        convert to standard form, enumerate all choices of two basic
        variables, determine all basic feasible solutions, match them to
        the corners of the feasible region, and find the optimum.
  \item How many candidate bases (choices of basic variables) are there at
        most for
        \begin{itemize}[leftmargin=1.8em,itemsep=0.2ex]
          \item[(a)] a standard-form system with $m=2$ equations and $n=4$
                variables,
          \item[(b)] an LP with $2$ decision variables and $3$ constraints
                of type $\le$ (after adding slack variables),
          \item[(c)] an LP with $3$ decision variables and $3$ constraints
                of type $\le$ (after adding slack variables)?
        \end{itemize}
        Explain why the number of corners of the feasible region can be
        strictly smaller than this bound.
  \item Convert the LP
        \[
          \max\; z=x_{1}+2x_{2}
          \quad\text{s.t.}\quad
          x_{1}+x_{2}\le 3,\;\;
          x_{1}\le 3,\;\;
          x_{1},x_{2}\ge 0
        \]
        to standard form and enumerate all basic solutions. Show that
        several bases produce the same corner, identify the degenerate
        basic solutions, and find the optimum.
  \item Convert to a maximisation in standard form:
        \[
          \min\; z=2x_{1}+3x_{2}-x_{3}
          \quad\text{s.t.}\quad
          \begin{aligned}[t]
            x_{1}+x_{2}+x_{3}&\le 10,\\
            x_{1}-x_{2}&\ge 2,\\
            x_{2}+x_{3}&=4,
          \end{aligned}
          \qquad x_{1},x_{2}\ge 0,\; x_{3}\ \text{free}.
        \]
\end{problems}

\section{Solutions}

\begin{solutions}
  \item Both constraints are of type $\le$, so one slack variable per row:
        \[
          \begin{aligned}
            \max\;\; & z=3x_{1}+2x_{2}+0s_{1}+0s_{2}\\
            \text{s.t.}\;\; & x_{1}+x_{2}+s_{1}=4\\
            & x_{1}+3x_{2}+s_{2}=6\\
            & x_{1},x_{2},s_{1},s_{2}\ge 0 .
          \end{aligned}
        \]
        $s_{1}$ and $s_{2}$ measure the unused capacity of the two
        constraints and earn nothing in the objective.

  \item The first constraint gets a slack variable ($+s_{1}$), the second
        is of type $\ge$ and gets a \emph{surplus} variable ($-s_{2}$):
        \[
          \begin{aligned}
            \max\;\; & z=2x_{1}+5x_{2}+0s_{1}+0s_{2}\\
            \text{s.t.}\;\; & x_{1}+2x_{2}+s_{1}=8\\
            & 3x_{1}+x_{2}-s_{2}=3\\
            & x_{1},x_{2},s_{1},s_{2}\ge 0 .
          \end{aligned}
        \]

  \item Two steps. \emph{Step 1 (min $\to$ max):} minimising
        $4x_{1}+3x_{2}$ is equivalent to maximising
        $w=-4x_{1}-3x_{2}$. \emph{Step 2 (surplus variables):} both
        constraints are of type $\ge$:
        \[
          \begin{aligned}
            \max\;\; & w=-4x_{1}-3x_{2}+0s_{1}+0s_{2}\\
            \text{s.t.}\;\; & 2x_{1}+x_{2}-s_{1}=10\\
            & x_{1}+3x_{2}-s_{2}=15\\
            & x_{1},x_{2},s_{1},s_{2}\ge 0 .
          \end{aligned}
        \]
        The optimal points are the same; the original minimum is recovered
        as $z^{*}=-w^{*}$.

  \item The equality constraint is copied unchanged; only the inequality
        gets a slack variable:
        \[
          \begin{aligned}
            \max\;\; & z=x_{1}+4x_{2}+0s_{1}\\
            \text{s.t.}\;\; & x_{1}+x_{2}=5\\
            & x_{1}+2x_{2}+s_{1}=8\\
            & x_{1},x_{2},s_{1}\ge 0 .
          \end{aligned}
        \]
        Note: the equation receives no new variable; adding one anyway is
        a common mistake.

  \item The free variable is split: $x_{2}=x_{2}^{+}-x_{2}^{-}$ with
        $x_{2}^{+},x_{2}^{-}\ge 0$. Substituting everywhere and adding
        slacks:
        \[
          \begin{aligned}
            \max\;\; & z=2x_{1}+x_{2}^{+}-x_{2}^{-}\\
            \text{s.t.}\;\; & x_{1}+x_{2}^{+}-x_{2}^{-}+s_{1}=3\\
            & x_{1}-x_{2}^{+}+x_{2}^{-}+s_{2}=2\\
            & x_{1},x_{2}^{+},x_{2}^{-},s_{1},s_{2}\ge 0 .
          \end{aligned}
        \]
        After solving, $x_{2}=x_{2}^{+}-x_{2}^{-}$ restores the original
        variable.

  \item Since $x_{2}\le 0$, substitute $y_{2}=-x_{2}\ge 0$, i.e.\
        $x_{2}=-y_{2}$. The objective becomes $3x_{1}-2(-y_{2})
        =3x_{1}+2y_{2}$; the constraints become $x_{1}+y_{2}\le 6$ and
        $2x_{1}-y_{2}\le 4$. With slacks:
        \[
          \begin{aligned}
            \max\;\; & z=3x_{1}+2y_{2}\\
            \text{s.t.}\;\; & x_{1}+y_{2}+s_{1}=6\\
            & 2x_{1}-y_{2}+s_{2}=4\\
            & x_{1},y_{2},s_{1},s_{2}\ge 0 .
          \end{aligned}
        \]
        At the end, $x_{2}=-y_{2}$.

  \item With
        \[
          c=\begin{pmatrix}5\\4\\3\end{pmatrix},\qquad
          x=\begin{pmatrix}x_{1}\\x_{2}\\x_{3}\end{pmatrix},\qquad
          A=\begin{bmatrix}
              2 & 3 & 1\\
              4 & 1 & 2\\
              3 & 4 & 2
            \end{bmatrix},\qquad
          b=\begin{pmatrix}5\\11\\8\end{pmatrix},
        \]
        the LP reads $\max\,c^{T}x$ s.t.\ $Ax\le b$, $x\ge 0$. Here $A$ is
        a $3\times 3$ matrix ($m=3$ constraints, $n=3$ variables),
        $b\in\R^{3}$, $c,x\in\R^{3}$.

  \item Standard form with slacks $s_{1},s_{2}$:
        \[
          \begin{aligned}
            \max\;\; & z=4x_{1}+3x_{2}+0s_{1}+0s_{2}\\
            \text{s.t.}\;\; & 2x_{1}+x_{2}+s_{1}=10\\
            & x_{1}+3x_{2}+s_{2}=15\\
            & x_{1},x_{2},s_{1},s_{2}\ge 0 .
          \end{aligned}
        \]
        In matrix notation $\bar{A}\bar{x}=b$, $\bar{x}\ge 0$ with
        \[
          \bar{A}=\begin{bmatrix}
              2 & 1 & 1 & 0\\
              1 & 3 & 0 & 1
            \end{bmatrix},\qquad
          \bar{x}=\begin{pmatrix}x_{1}\\x_{2}\\s_{1}\\s_{2}\end{pmatrix},
          \qquad
          b=\begin{pmatrix}10\\15\end{pmatrix},\qquad
          \bar{c}=\begin{pmatrix}4\\3\\0\\0\end{pmatrix}.
        \]
        The slack columns form an identity block, and the slack costs are
        zero: this identity block will serve as the starting basis of
        the simplex method.

  \item Here $m=2$, $n=4$, so there are $\binom{4}{2}=6$ candidate pairs.
        Set the other two variables to $0$ each time and solve.
        \begin{itemize}[leftmargin=1.8em,itemsep=0.3ex]
          \item $\{x_{1},x_{2}\}$: $x_{3}=x_{4}=0$. Row~2: $x_{2}=2$;
                row~1: $x_{1}=4-2=2$. Solution $(2,2,0,0)$, feasible.
          \item $\{x_{1},x_{3}\}$: the columns of $x_{1}$ and $x_{3}$ are
                $\begin{pmatrix}1\\0\end{pmatrix}$ and
                $\begin{pmatrix}1\\0\end{pmatrix}$: linearly dependent,
                \emph{no basis} (indeed $x_{2}=x_{4}=0$ makes row~2 read
                $0=2$).
          \item $\{x_{1},x_{4}\}$: $x_{2}=x_{3}=0$. Row~1: $x_{1}=4$;
                row~2: $x_{4}=2$. Solution $(4,0,0,2)$, feasible.
          \item $\{x_{2},x_{3}\}$: $x_{1}=x_{4}=0$. Row~2: $x_{2}=2$;
                row~1: $x_{3}=4-2=2$. Solution $(0,2,2,0)$, feasible.
          \item $\{x_{2},x_{4}\}$: $x_{1}=x_{3}=0$. Row~1: $x_{2}=4$;
                row~2: $4+x_{4}=2$, so $x_{4}=-2$. Solution $(0,4,0,-2)$,
                basic but \emph{infeasible}.
          \item $\{x_{3},x_{4}\}$: $x_{1}=x_{2}=0$. Row~1: $x_{3}=4$;
                row~2: $x_{4}=2$. Solution $(0,0,4,2)$, feasible.
        \end{itemize}
        Summary: $6$ candidate pairs, $5$ bases, $4$ basic feasible
        solutions, $1$ infeasible basic solution.

  \item Check the two equations, count the nonzero entries (a basic
        solution has at most $m=2$), and check the independence of the
        corresponding columns.
        \begin{itemize}[leftmargin=1.8em,itemsep=0.3ex]
          \item[(a)] $(2,2,0,0)$: row~1: $2+2+0=4$~\checkmark; row~2:
                $4+2+0=6$~\checkmark. Nonzero entries $x_{1},x_{2}$ with
                columns
                $\begin{pmatrix}1\\2\end{pmatrix},
                 \begin{pmatrix}1\\1\end{pmatrix}$
                (determinant $1\cdot 1-1\cdot 2=-1\ne 0$): a
                \emph{basic feasible solution}.
          \item[(b)] $(1,2,1,2)$: row~1: $1+2+1=4$~\checkmark; row~2:
                $2+2+2=6$~\checkmark. Feasible, but four positive entries
                $>m=2$: \emph{feasible, not basic}.
          \item[(c)] $(3,0,1,0)$: row~1: $3+0+1=4$~\checkmark; row~2:
                $6+0+0=6$~\checkmark. Nonzero entries $x_{1},x_{3}$ with
                columns
                $\begin{pmatrix}1\\2\end{pmatrix},
                 \begin{pmatrix}1\\0\end{pmatrix}$
                (determinant $-2\ne 0$): a \emph{basic feasible solution}.
          \item[(d)] $(0,6,-2,0)$: row~1: $0+6-2=4$~\checkmark; row~2:
                $0+6+0=6$~\checkmark. A solution, and basic (columns of
                $x_{2},x_{3}$ are independent), but $x_{3}=-2<0$:
                \emph{basic, not feasible}.
          \item[(e)] $(1,1,1,1)$: row~1: $1+1+1=3\ne 4$, \emph{not even
                a solution} of the system.
        \end{itemize}

  \item Standard form (from Problem~1):
        \[
          x_{1}+x_{2}+s_{1}=4,\qquad
          x_{1}+3x_{2}+s_{2}=6,\qquad
          x_{1},x_{2},s_{1},s_{2}\ge 0 .
        \]
        All $\binom{4}{2}=6$ pairs of columns are linearly independent
        here, so each pair gives exactly one basic solution:
        \begin{center}
        \begin{tabular}{llll}
          \hline
          basic variables & nonbasic $=0$ &
          $(x_{1},x_{2},s_{1},s_{2})$ & feasible?\\
          \hline
          $\{s_{1},s_{2}\}$ & $x_{1},x_{2}$ & $(0,0,4,6)$ & yes (BFS)\\
          $\{x_{1},s_{2}\}$ & $x_{2},s_{1}$ & $(4,0,0,2)$ & yes (BFS)\\
          $\{x_{1},s_{1}\}$ & $x_{2},s_{2}$ & $(6,0,-2,0)$ & no\\
          $\{x_{2},s_{1}\}$ & $x_{1},s_{2}$ & $(0,2,2,0)$ & yes (BFS)\\
          $\{x_{2},s_{2}\}$ & $x_{1},s_{1}$ & $(0,4,0,-6)$ & no\\
          $\{x_{1},x_{2}\}$ & $s_{1},s_{2}$ & $(3,1,0,0)$ & yes (BFS)\\
          \hline
        \end{tabular}
        \end{center}
        Sample computations: for $\{x_{1},s_{1}\}$, row~2 with
        $x_{2}=s_{2}=0$ gives $x_{1}=6$, then row~1 gives
        $s_{1}=4-6=-2<0$. For $\{x_{1},x_{2}\}$, subtracting the equations
        $x_{1}+x_{2}=4$ and $x_{1}+3x_{2}=6$ gives $2x_{2}=2$, so
        $x_{2}=1$, $x_{1}=3$. There are $4$ basic feasible solutions.

  \item Corners are intersections of two active constraint lines that
        satisfy all constraints:
        \begin{itemize}[leftmargin=1.8em,itemsep=0.3ex]
          \item $x_{1}=0$, $x_{2}=0$: $(0,0)$, feasible.
          \item $x_{2}=0$, $x_{1}+x_{2}=4$: $(4,0)$; check
                $4+0=4\le 6$~\checkmark, corner.
          \item $x_{2}=0$, $x_{1}+3x_{2}=6$: $(6,0)$; but $6>4$ violates
                the first constraint, \emph{not} a corner.
          \item $x_{1}=0$, $x_{1}+3x_{2}=6$: $(0,2)$; check
                $0+2=2\le 4$~\checkmark, corner.
          \item $x_{1}=0$, $x_{1}+x_{2}=4$: $(0,4)$; but
                $0+12=12>6$, \emph{not} a corner.
          \item $x_{1}+x_{2}=4$, $x_{1}+3x_{2}=6$: subtracting,
                $2x_{2}=2$, so $(3,1)$; feasible, corner.
        \end{itemize}
        Matching with Problem~11 ($z=3x_{1}+2x_{2}$):
        \begin{center}
        \begin{tabular}{llll}
          \hline
          corner $(x_{1},x_{2})$ & BFS $(x_{1},x_{2},s_{1},s_{2})$ &
          basis & $z$\\
          \hline
          $(0,0)$ & $(0,0,4,6)$ & $\{s_{1},s_{2}\}$ & $0$\\
          $(4,0)$ & $(4,0,0,2)$ & $\{x_{1},s_{2}\}$ & $12$\\
          $(0,2)$ & $(0,2,2,0)$ & $\{x_{2},s_{1}\}$ & $4$\\
          $(3,1)$ & $(3,1,0,0)$ & $\{x_{1},x_{2}\}$ & $11$\\
          \hline
        \end{tabular}
        \end{center}
        The two infeasible basic solutions are precisely the two rejected
        intersection points: $(6,0,-2,0)$ sits at $(6,0)$ and
        $(0,4,0,-6)$ at $(0,4)$, intersections of constraint lines
        \emph{outside} the feasible region. The maximum of $z$ over the
        BFS is $z^{*}=12$ at $(x_{1},x_{2})=(4,0)$, where $s_{1}=0$ (first
        constraint binding) and $s_{2}=2$ (second constraint slack).

  \item Standard form:
        \[
          x_{1}+2x_{2}+s_{1}=6,\qquad
          x_{1}+s_{2}=4,\qquad
          x_{1},x_{2},s_{1},s_{2}\ge 0 .
        \]
        Going through the $\binom{4}{2}=6$ pairs:
        \begin{itemize}[leftmargin=1.8em,itemsep=0.3ex]
          \item $\{s_{1},s_{2}\}$: $(0,0,6,4)$, BFS, corner $(0,0)$.
          \item $\{x_{1},s_{1}\}$: row~2 gives $x_{1}=4$, row~1 gives
                $s_{1}=6-4=2$: $(4,0,2,0)$, BFS, corner $(4,0)$.
          \item $\{x_{1},s_{2}\}$: row~1 gives $x_{1}=6$, row~2 gives
                $s_{2}=4-6=-2$: $(6,0,0,-2)$: infeasible (the point
                $(6,0)$ lies outside).
          \item $\{x_{2},s_{1}\}$: columns
                $\begin{pmatrix}2\\0\end{pmatrix}$ and
                $\begin{pmatrix}1\\0\end{pmatrix}$ are dependent:
                \emph{no basis} (row~2 would read $0=4$).
          \item $\{x_{2},s_{2}\}$: row~1 gives $2x_{2}=6$, so $x_{2}=3$;
                row~2 gives $s_{2}=4$: $(0,3,0,4)$, BFS, corner
                $(0,3)$.
          \item $\{x_{1},x_{2}\}$: row~2 gives $x_{1}=4$; row~1 gives
                $4+2x_{2}=6$, so $x_{2}=1$: $(4,1,0,0)$, BFS, corner
                $(4,1)$.
        \end{itemize}
        Objective values $z=2x_{1}+3x_{2}$ at the four corners:
        \[
          z(0,0)=0,\qquad z(4,0)=8,\qquad z(0,3)=9,\qquad z(4,1)=11 .
        \]
        Optimum: $z^{*}=11$ at $(x_{1},x_{2})=(4,1)$, basis
        $\{x_{1},x_{2}\}$.

  \item Counting with the binomial coefficient:
        \begin{itemize}[leftmargin=1.8em,itemsep=0.3ex]
          \item[(a)] $\dbinom{4}{2}=\dfrac{4!}{2!\,2!}=6$.
          \item[(b)] Slacks make $n=2+3=5$ variables and $m=3$ equations:
                $\dbinom{5}{3}=10$.
          \item[(c)] $n=3+3=6$, $m=3$:
                $\dbinom{6}{3}=\dfrac{6\cdot 5\cdot 4}{3\cdot 2\cdot 1}
                =20$.
        \end{itemize}
        The bound overcounts for three reasons: a chosen column set can be
        linearly dependent (no basic solution at all), a basic solution
        can have negative entries (an intersection point outside the
        feasible region), and in the degenerate case several bases
        describe the same corner. Hence: number of corners
        $\le\binom{n}{m}$, usually with strict inequality.

  \item Standard form:
        \[
          x_{1}+x_{2}+s_{1}=3,\qquad
          x_{1}+s_{2}=3,\qquad
          x_{1},x_{2},s_{1},s_{2}\ge 0 .
        \]
        The six pairs:
        \begin{itemize}[leftmargin=1.8em,itemsep=0.3ex]
          \item $\{s_{1},s_{2}\}$: $(0,0,3,3)$, BFS, corner $(0,0)$.
          \item $\{x_{2},s_{2}\}$: row~1: $x_{2}=3$; row~2: $s_{2}=3$:
                $(0,3,0,3)$, BFS, corner $(0,3)$.
          \item $\{x_{2},s_{1}\}$: columns
                $\begin{pmatrix}1\\0\end{pmatrix}$,
                $\begin{pmatrix}1\\0\end{pmatrix}$ dependent, no basis.
          \item $\{x_{1},x_{2}\}$: row~2: $x_{1}=3$; row~1: $x_{2}=0$:
                $(3,0,0,0)$, BFS, \emph{degenerate} (basic variable
                $x_{2}=0$).
          \item $\{x_{1},s_{1}\}$: row~2: $x_{1}=3$; row~1: $s_{1}=0$:
                $(3,0,0,0)$, the \emph{same} point, again degenerate.
          \item $\{x_{1},s_{2}\}$: row~1: $x_{1}=3$; row~2: $s_{2}=0$:
                $(3,0,0,0)$, the same point a third time.
        \end{itemize}
        Three different bases, $\{x_{1},x_{2}\}$, $\{x_{1},s_{1}\}$,
        $\{x_{1},s_{2}\}$, all represent the corner $(3,0)$: there, the
        three lines $x_{2}=0$, $x_{1}+x_{2}=3$ and $x_{1}=3$ all meet
        (the constraint $x_{1}\le 3$ is redundant and touches the region
        only in this point). The region has just $3$ corners: $(0,0)$,
        $(3,0)$, $(0,3)$. Objective $z=x_{1}+2x_{2}$:
        \[
          z(0,0)=0,\qquad z(3,0)=3,\qquad z(0,3)=6 ,
        \]
        so $z^{*}=6$ at $(0,3)$.

  \item Three conversions at once. \emph{Min $\to$ max:} maximise
        $w=-2x_{1}-3x_{2}+x_{3}$. \emph{Free variable:}
        $x_{3}=u-v$, $u,v\ge 0$. \emph{Constraints:} slack for the
        $\le$ row, surplus for the $\ge$ row, the equation stays:
        \[
          \begin{aligned}
            \max\;\; & w=-2x_{1}-3x_{2}+u-v\\
            \text{s.t.}\;\; & x_{1}+x_{2}+u-v+s_{1}=10\\
            & x_{1}-x_{2}-s_{2}=2\\
            & x_{2}+u-v=4\\
            & x_{1},x_{2},u,v,s_{1},s_{2}\ge 0 .
          \end{aligned}
        \]
        Afterwards $x_{3}=u-v$ and $z^{*}=-w^{*}$.
\end{solutions}

\section{Additional Problems}

\begin{problems}
  \item Convert to standard form (watch the last constraint):
        \[
          \max\; z=4x_{1}+x_{2}
          \quad\text{s.t.}\quad
          x_{1}+2x_{2}\le 6,\;\;
          3x_{1}+2x_{2}\ge 6,\;\;
          x_{1}-x_{2}\ge -2,\;\;
          x_{1},x_{2}\ge 0 .
        \]
  \item Convert to a maximisation in standard form:
        \[
          \min\; z=x_{1}-2x_{2}+x_{3}
          \quad\text{s.t.}\quad
          x_{1}+x_{2}-x_{3}=4,\;\;
          x_{1}-x_{2}\le 2,\;\;
          x_{1},x_{2}\ge 0,\; x_{3}\ \text{free}.
        \]
  \item Write the LP
        \[
          \max\; z=2x_{1}+3x_{2}+x_{3}
          \quad\text{s.t.}\quad
          x_{1}+x_{2}+x_{3}\le 6,\;\;
          x_{1}+2x_{2}\le 4,\;\;
          x_{1},x_{2},x_{3}\ge 0
        \]
        in the matrix form $\max\,c^{T}x$ s.t.\ $Ax\le b$, $x\ge 0$, then
        give the standard-form matrix $\bar{A}$ and the extended cost
        vector $\bar{c}$. How many candidate bases does the standard form
        have?
  \item Find all basic solutions of
        \[
          \begin{aligned}
            2x_{1}+x_{2}+x_{3} &= 4\\
            x_{1}+x_{2}+x_{4} &= 3
          \end{aligned}
          \qquad x_{1},\dots,x_{4}\ge 0 ,
        \]
        and state which are feasible.
  \item For the LP
        \[
          \max\; z=x_{1}+x_{2}
          \quad\text{s.t.}\quad
          2x_{1}+x_{2}\le 6,\;\;
          x_{1}+2x_{2}\le 6,\;\;
          x_{1},x_{2}\ge 0 ,
        \]
        enumerate all basic solutions of the standard form, match the
        basic feasible solutions to the corners of the feasible region,
        and determine the optimum.
  \item Consider the system
        \[
          \begin{aligned}
            x_{1}+x_{2}+x_{3} &= 3\\
            x_{2}+x_{4} &= 1
          \end{aligned}
          \qquad x_{1},\dots,x_{4}\ge 0 .
        \]
        Classify each vector as basic feasible / feasible but not basic /
        basic but not feasible / not a solution:
        \begin{itemize}[leftmargin=1.8em,itemsep=0.2ex]
          \item[(a)] $(2,1,0,0)$
          \item[(b)] $(1,1,1,0)$
          \item[(c)] $(3,0,0,1)$
          \item[(d)] $(0,3,0,-2)$
        \end{itemize}
  \item How many candidate bases are there at most for
        \begin{itemize}[leftmargin=1.8em,itemsep=0.2ex]
          \item[(a)] a standard-form system with $m=2$ equations and
                $n=5$ variables,
          \item[(b)] an LP with $4$ decision variables and $2$ constraints
                of type $\le$,
          \item[(c)] an LP with $2$ decision variables and $4$ constraints
                of type $\le$?
        \end{itemize}
  \item For the LP
        \[
          \max\; z=x_{1}+x_{2}
          \quad\text{s.t.}\quad
          x_{1}\le 2,\;\;
          x_{2}\le 2,\;\;
          x_{1}+x_{2}\le 4,\;\;
          x_{1},x_{2}\ge 0 ,
        \]
        write the standard form ($m=3$, $n=5$) and show by direct
        computation that the three bases $\{x_{1},x_{2},s_{1}\}$,
        $\{x_{1},x_{2},s_{2}\}$ and $\{x_{1},x_{2},s_{3}\}$ all yield the
        same basic solution. Which corner is it, is it degenerate, and
        what is the optimal value?
\end{problems}

\section{Additional Solutions}

\begin{solutions}
  \item The third constraint has a negative right-hand side; multiply it
        by $-1$ first: $x_{1}-x_{2}\ge -2$ becomes $-x_{1}+x_{2}\le 2$.
        Then slack, surplus, slack:
        \[
          \begin{aligned}
            \max\;\; & z=4x_{1}+x_{2}\\
            \text{s.t.}\;\; & x_{1}+2x_{2}+s_{1}=6\\
            & 3x_{1}+2x_{2}-s_{2}=6\\
            & -x_{1}+x_{2}+s_{3}=2\\
            & x_{1},x_{2},s_{1},s_{2},s_{3}\ge 0 .
          \end{aligned}
        \]

  \item Maximise $w=-x_{1}+2x_{2}-x_{3}$ and split $x_{3}=u-v$ with
        $u,v\ge 0$ (so $-x_{3}=-u+v$). The equation stays, the inequality
        gets a slack:
        \[
          \begin{aligned}
            \max\;\; & w=-x_{1}+2x_{2}-u+v\\
            \text{s.t.}\;\; & x_{1}+x_{2}-u+v=4\\
            & x_{1}-x_{2}+s_{1}=2\\
            & x_{1},x_{2},u,v,s_{1}\ge 0 .
          \end{aligned}
        \]
        Back-substitution: $x_{3}=u-v$ and $z^{*}=-w^{*}$.

  \item Matrix (inequality) form:
        \[
          c=\begin{pmatrix}2\\3\\1\end{pmatrix},\qquad
          A=\begin{bmatrix}
              1 & 1 & 1\\
              1 & 2 & 0
            \end{bmatrix},\qquad
          b=\begin{pmatrix}6\\4\end{pmatrix}.
        \]
        Standard form with slacks $s_{1},s_{2}$:
        \[
          \bar{A}=\begin{bmatrix}
              1 & 1 & 1 & 1 & 0\\
              1 & 2 & 0 & 0 & 1
            \end{bmatrix},\qquad
          \bar{x}=\begin{pmatrix}x_{1}\\x_{2}\\x_{3}\\s_{1}\\s_{2}
            \end{pmatrix},\qquad
          \bar{c}=\begin{pmatrix}2\\3\\1\\0\\0\end{pmatrix}.
        \]
        With $n=5$ variables and $m=2$ equations there are at most
        $\binom{5}{2}=10$ candidate bases.

  \item All $\binom{4}{2}=6$ column pairs are linearly independent here,
        so each yields a basic solution:
        \begin{center}
        \begin{tabular}{llll}
          \hline
          basic variables & nonbasic $=0$ &
          $(x_{1},x_{2},x_{3},x_{4})$ & feasible?\\
          \hline
          $\{x_{1},x_{2}\}$ & $x_{3},x_{4}$ & $(1,2,0,0)$ & yes (BFS)\\
          $\{x_{1},x_{3}\}$ & $x_{2},x_{4}$ & $(3,0,-2,0)$ & no\\
          $\{x_{1},x_{4}\}$ & $x_{2},x_{3}$ & $(2,0,0,1)$ & yes (BFS)\\
          $\{x_{2},x_{3}\}$ & $x_{1},x_{4}$ & $(0,3,1,0)$ & yes (BFS)\\
          $\{x_{2},x_{4}\}$ & $x_{1},x_{3}$ & $(0,4,0,-1)$ & no\\
          $\{x_{3},x_{4}\}$ & $x_{1},x_{2}$ & $(0,0,4,3)$ & yes (BFS)\\
          \hline
        \end{tabular}
        \end{center}
        Sample computations: $\{x_{1},x_{2}\}$: subtracting
        $x_{1}+x_{2}=3$ from $2x_{1}+x_{2}=4$ gives $x_{1}=1$, so
        $x_{2}=2$. $\{x_{1},x_{3}\}$: row~2 gives $x_{1}=3$, row~1 gives
        $x_{3}=4-6=-2<0$. $\{x_{2},x_{4}\}$: row~1 gives $x_{2}=4$,
        row~2 gives $x_{4}=3-4=-1<0$. Four of the six basic solutions are
        feasible.

  \item Standard form:
        \[
          2x_{1}+x_{2}+s_{1}=6,\qquad
          x_{1}+2x_{2}+s_{2}=6,\qquad
          x_{1},x_{2},s_{1},s_{2}\ge 0 .
        \]
        All six pairs are bases:
        \begin{center}
        \begin{tabular}{lllll}
          \hline
          basis & $(x_{1},x_{2},s_{1},s_{2})$ & feasible? &
          point & $z=x_{1}+x_{2}$\\
          \hline
          $\{s_{1},s_{2}\}$ & $(0,0,6,6)$ & yes & $(0,0)$ & $0$\\
          $\{x_{1},s_{2}\}$ & $(3,0,0,3)$ & yes & $(3,0)$ & $3$\\
          $\{x_{1},s_{1}\}$ & $(6,0,-6,0)$ & no & & \\
          $\{x_{2},s_{1}\}$ & $(0,3,3,0)$ & yes & $(0,3)$ & $3$\\
          $\{x_{2},s_{2}\}$ & $(0,6,0,-6)$ & no & & \\
          $\{x_{1},x_{2}\}$ & $(2,2,0,0)$ & yes & $(2,2)$ & $4$\\
          \hline
        \end{tabular}
        \end{center}
        Sample computations: $\{x_{1},x_{2}\}$: doubling
        $2x_{1}+x_{2}=6$ gives $4x_{1}+2x_{2}=12$; subtracting
        $x_{1}+2x_{2}=6$ gives $3x_{1}=6$, so $x_{1}=2$, $x_{2}=2$.
        $\{x_{1},s_{1}\}$: row~2 gives $x_{1}=6$, row~1 gives
        $s_{1}=6-12=-6$. The corners are $(0,0)$, $(3,0)$, $(2,2)$,
        $(0,3)$; the two infeasible basic solutions sit at the outside
        intersections $(6,0)$ and $(0,6)$. Optimum: $z^{*}=4$ at
        $(x_{1},x_{2})=(2,2)$.

  \item With $m=2$, a basic solution has at most two nonzero entries.
        \begin{itemize}[leftmargin=1.8em,itemsep=0.3ex]
          \item[(a)] $(2,1,0,0)$: row~1: $2+1+0=3$~\checkmark; row~2:
                $1+0=1$~\checkmark. Columns of $x_{1},x_{2}$:
                $\begin{pmatrix}1\\0\end{pmatrix},
                 \begin{pmatrix}1\\1\end{pmatrix}$, independent,
                \emph{basic feasible}.
          \item[(b)] $(1,1,1,0)$: row~1: $3$~\checkmark; row~2:
                $1$~\checkmark. Three positive entries $>2$,
                \emph{feasible, not basic}.
          \item[(c)] $(3,0,0,1)$: row~1: $3$~\checkmark; row~2:
                $0+1=1$~\checkmark. Columns of $x_{1},x_{4}$:
                $\begin{pmatrix}1\\0\end{pmatrix},
                 \begin{pmatrix}0\\1\end{pmatrix}$, independent,
                \emph{basic feasible}.
          \item[(d)] $(0,3,0,-2)$: row~1: $0+3+0=3$~\checkmark; row~2:
                $3-2=1$~\checkmark. Columns of $x_{2},x_{4}$ independent,
                but $x_{4}=-2<0$, \emph{basic, not feasible}.
        \end{itemize}

  \item Counting:
        \begin{itemize}[leftmargin=1.8em,itemsep=0.3ex]
          \item[(a)] $\dbinom{5}{2}=10$.
          \item[(b)] $n=4+2=6$, $m=2$: $\dbinom{6}{2}=15$.
          \item[(c)] $n=2+4=6$, $m=4$: $\dbinom{6}{4}=\dbinom{6}{2}=15$.
        \end{itemize}
        Parts (b) and (c) coincide because
        $\binom{6}{4}=\binom{6}{2}$: choosing the $4$ basic variables is
        the same as choosing the $2$ nonbasic ones.

  \item Standard form ($m=3$ equations, $n=5$ variables):
        \[
          x_{1}+s_{1}=2,\qquad
          x_{2}+s_{2}=2,\qquad
          x_{1}+x_{2}+s_{3}=4,\qquad
          \text{all}\ \ge 0 .
        \]
        \begin{itemize}[leftmargin=1.8em,itemsep=0.3ex]
          \item Basis $\{x_{1},x_{2},s_{1}\}$, nonbasic $s_{2}=s_{3}=0$:
                row~2 gives $x_{2}=2$; row~3 gives $x_{1}=4-2=2$; row~1
                gives $s_{1}=2-2=0$. Solution $(2,2,0,0,0)$.
          \item Basis $\{x_{1},x_{2},s_{2}\}$, nonbasic $s_{1}=s_{3}=0$:
                row~1 gives $x_{1}=2$; row~3 gives $x_{2}=4-2=2$; row~2
                gives $s_{2}=2-2=0$. Solution $(2,2,0,0,0)$.
          \item Basis $\{x_{1},x_{2},s_{3}\}$, nonbasic $s_{1}=s_{2}=0$:
                row~1 gives $x_{1}=2$; row~2 gives $x_{2}=2$; row~3 gives
                $s_{3}=4-4=0$. Solution $(2,2,0,0,0)$.
        \end{itemize}
        All three bases yield the corner $(x_{1},x_{2})=(2,2)$, and in
        each case one \emph{basic} variable equals $0$ (only $2<m=3$
        nonzero entries): the basic solution is \emph{degenerate}.
        Geometrically, the three lines $x_{1}=2$, $x_{2}=2$ and
        $x_{1}+x_{2}=4$ all pass through $(2,2)$. The optimal value is
        $z^{*}=2+2=4$, attained exactly at this degenerate corner.
\end{solutions}
